\section{Algebraic Effects in OCaml}
This section provides background on effects and handlers; readers focused on
\tempo{} can skip to the overview and return here for implementation details.

\paragraph{One-sentence intuition}
An effect captures "what to do next" (the continuation), and the scheduler
decides \emph{when} to resume it, which aligns naturally with instant-based
reactive execution.
Algebraic effects provide a structured way to model effectful operations by
separating the declaration of an effect from its interpretation. Instead of
hardwiring side effects into the runtime or encoding them as monads, a program
\emph{performs} an effect at a call site, and a handler decides how to interpret
it. The handler can resume the captured continuation immediately, defer it, or
even resume it multiple times, which makes effects a natural fit for
scheduling, suspension, and user-defined control operators.

In OCaml 5, effect handlers are first-class and integrate with the language
runtime. The model follows the theory of algebraic effects and handlers
\cite{plotkinpretnar2009}, and it has been validated by practical systems such
as Eff \cite{eff2013} and the Multicore OCaml project
\cite{multicoreocaml2020}. For our purposes, the important point is that the
handler sees the point of suspension as a value---a \emph{delimited
continuation}---and can choose precisely when to resume it. Algebraic effects
can therefore be viewed as a disciplined interface to delimited continuations,
with the handler providing the delimiter \cite{danvyfilinski1990,filinski1994}.

From the programmer's perspective, effects look like ordinary function calls
with additional control semantics. For instance, an \texttt{Await} effect can be
implemented so that it captures the continuation and stores it in a signal's
waiter list, rather than resuming it immediately. A \texttt{Pause} effect can
package the continuation into a task and enqueue it for the next instant.
Similarly, \texttt{Emit} can be interpreted as a state update plus a resumption
of waiting tasks. The key idea is that the operational behavior is defined in
the handler, not at the call site.

This style is particularly well suited to synchronous reactive runtimes, where
the semantics depend on \emph{when} computations resume. By choosing whether a
continuation is resumed in the current instant or deferred to the next one, the
handler can implement the delayed-absence discipline and maintain determinism.
In other words, algebraic effects provide the control lever that makes an
instant-based scheduler possible at library level.

For \tempo{}, algebraic effects are therefore not an implementation detail but the
core enabler: they allow the library to encode reactive constructs (signals,
guards, preemption) directly as effectful operations. The scheduler lives in
user space, with explicit queues and signal registries, and the effect handler
defines the boundary between instantaneous reactions and deferred computation.

\paragraph{Deep vs shallow handlers}
OCaml 5 distinguishes \emph{deep} handlers, where the continuation is
automatically re-wrapped by the same handler upon resumption, from
\emph{shallow} handlers, where the resuming code runs outside the handler
unless it is re-installed explicitly. \tempo{} uses deep handlers (via
\texttt{Effect.Deep}) to ensure that every resumed task continues under the
same scheduler, which is essential for consistent instant semantics.

\paragraph{Cooperative threads via effects}
It is well known that algebraic effects and delimited continuations can be
used to implement cooperative threading. A \texttt{Yield} or \texttt{Pause}
effect captures the current continuation and enqueues it, while the scheduler
selects another task to run. This is the same pattern used in many effect-based
thread libraries and serves as the foundation for \tempo{}'s instant-based
cooperative scheduler.

\paragraph{Intuition (small example)}
The following toy sketch shows how an effect is performed and how a handler
decides when to resume the continuation.
\begin{lstlisting}[style=tempo]
(* Effect declaration *)
type _ Effect.t +=
  | Yield : unit Effect.t

let yield () = perform Yield

(* Handler decides when to resume *)
let run f =
  match f () with
  | () -> ()
  | effect Yield k ->
      (* Defer continuation; resume later *)
      Queue.add (fun () -> continue k ()) pending
\end{lstlisting}

\begin{figure}[h]
\centering
\resizebox{\columnwidth}{!}{%
\begin{tikzpicture}[
  node distance=10mm,
  box/.style=tempo,
  draw,
  rounded corners,
  align=center,
  minimum width=34mm,
  minimum height=8mm,
  arrow/.style=tempo
]
\node[box] (effect) {Effect performed\\(Emit/Await/Pause/...)};
\node[box, below=of effect] (handler) {Effect handler\\captures continuation};
\node[box, below=of handler] (decision) {Schedule decision\\now or next instant};
\node[box, left=of decision] (now) {Enqueue\\current instant};
\node[box, right=of decision] (next) {Enqueue\\next instant};
\node[box, below=of decision] (resume) {Continuation resumes\\when dequeued};

\draw[arrow] (effect) -- (handler);
\draw[arrow] (handler) -- (decision);
\draw[arrow] (decision) -- (now);
\draw[arrow] (decision) -- (next);
\draw[arrow] (now) |- (resume);
\draw[arrow] (next) |- (resume);
\end{tikzpicture}%
}
\caption{Effect handling and continuation scheduling.}
\end{figure}

In the next section, we show how \tempo{} uses these mechanisms to implement a
synchronous runtime with explicit instants, signals, and control primitives.
